\documentclass[12pt,a4paper,openany,oneside]{book}

% --- Paquetes ---
\usepackage[utf8]{inputenc}
\usepackage[spanish, es-noshorthands]{babel}
\usepackage{amsmath, amsfonts, amssymb}
\usepackage{graphicx}
\usepackage{geometry}
\usepackage{hyperref}
\usepackage{booktabs}
\usepackage{fancyhdr}
\usepackage{listings}
\usepackage{xcolor}
\usepackage{tikz}
\usepackage{pgfplots}
\usetikzlibrary{arrows.meta, positioning, shapes.geometric}
\pgfplotsset{compat=1.18}

\geometry{margin=2.5cm}

% --- Colores y estilos para código Python ---
\definecolor{codegreen}{rgb}{0,0.6,0}
\definecolor{codegray}{rgb}{0.5,0.5,0.5}
\definecolor{codepurple}{rgb}{0.58,0,0.82}
\definecolor{backcolour}{rgb}{0.95,0.95,0.92}

\lstset{
    language=Python,
    backgroundcolor=\color{backcolour},   
    commentstyle=\color{codegreen},
    keywordstyle=\color{blue},
    numberstyle=\tiny\color{codegray},
    stringstyle=\color{codepurple},
    basicstyle=\ttfamily\small,
    breaklines=true,
    frame=single,
    showstringspaces=false
}

% --- Estilo de cabecera y pie ---
\pagestyle{fancy}
\fancyhf{}
\renewcommand{\headrulewidth}{0.4pt}
\renewcommand{\footrulewidth}{0.4pt}
\fancyhead[L]{\small Carlos César Sánchez Coronel}
\fancyhead[R]{\small Sesión 10: Espacios Vectoriales y Transformaciones Lineales}
\fancyfoot[L]{\small Carlos César Sánchez Coronel}
\fancyfoot[C]{\thepage}

% --- Portada ---
\title{
    \textbf{Sesión 10} \\ 
    \Huge \textbf{ESPACIOS VECTORIALES, CAMBIO DE BASE Y TRANSFORMACIONES LINEALES} \\
    \large La geometría de los datos y las redes neuronales \\
    \normalsize Aplicaciones en deep learning, NLP y visión por computador
}
\author{
    \small Por \\ \vspace{0.2cm}
    \Large \textsc{Carlos César Sánchez Coronel}
}
\date{2026}

\begin{document}

\maketitle
\tableofcontents
\addcontentsline{toc}{chapter}{Índice general}

% ------------------------------------------------------------------
% CAPÍTULO 10: SESIÓN 10 - ESPACIOS VECTORIALES Y TRANSFORMACIONES LINEALES
% ------------------------------------------------------------------
\chapter{Espacios Vectoriales, Cambio de Base y Transformaciones Lineales}

En esta sesión nos adentramos en el corazón del álgebra lineal: los espacios vectoriales y las transformaciones lineales. Estos conceptos no solo son bellos desde el punto de vista matemático, sino que son la base de cómo las redes neuronales procesan la información. Una capa fully connected no es más que una transformación lineal seguida de una no linealidad. El cambio de base nos permite entender los embeddings en NLP. Y las transformaciones afines (rotaciones, escalados) son esenciales en data augmentation para visión por computador. Al finalizar, serás capaz de visualizar cómo una matriz transforma un espacio y conectar estas ideas con aplicaciones reales en inteligencia artificial.

\section{Espacios vectoriales: el hogar de los datos}

\subsection{Definición intuitiva}
Un espacio vectorial es un conjunto de objetos (llamados vectores) que se pueden sumar entre sí y multiplicar por escalares (números reales), cumpliendo ciertas propiedades (asociatividad, conmutatividad, existencia de neutro, etc.). Los ejemplos más familiares son $\mathbb{R}^2$ (el plano) y $\mathbb{R}^3$ (el espacio tridimensional), pero también hay espacios de funciones, matrices, etc.

En inteligencia artificial, los datos viven en espacios vectoriales: cada punto es un vector de características, y las operaciones que realizamos (sumar, promediar, etc.) tienen sentido dentro de ese espacio.

\subsection{Definición formal}
Un espacio vectorial sobre $\mathbb{R}$ es un conjunto $V$ con dos operaciones:
\begin{itemize}
    \item Suma: $+ : V \times V \to V$, que satisface:
    \begin{align*}
        &\mathbf{u} + \mathbf{v} = \mathbf{v} + \mathbf{u} \quad \text{(conmutativa)}\\
        &(\mathbf{u} + \mathbf{v}) + \mathbf{w} = \mathbf{u} + (\mathbf{v} + \mathbf{w}) \quad \text{(asociativa)}\\
        &\exists \mathbf{0} \in V \text{ tal que } \mathbf{v} + \mathbf{0} = \mathbf{v} \quad \text{(neutro)}\\
        &\forall \mathbf{v} \in V, \exists (-\mathbf{v}) \in V \text{ tal que } \mathbf{v} + (-\mathbf{v}) = \mathbf{0} \quad \text{(inverso)}
    \end{align*}
    \item Producto por escalar: $\cdot : \mathbb{R} \times V \to V$, que satisface:
    \begin{align*}
        &a(\mathbf{v} + \mathbf{w}) = a\mathbf{v} + a\mathbf{w}\\
        &(a+b)\mathbf{v} = a\mathbf{v} + b\mathbf{v}\\
        &(ab)\mathbf{v} = a(b\mathbf{v})\\
        &1 \cdot \mathbf{v} = \mathbf{v}
    \end{align*}
\end{itemize}

\subsection{Subespacios}
Un subespacio vectorial es un subconjunto $W \subseteq V$ que por sí mismo es un espacio vectorial (es decir, cerrado bajo suma y producto por escalar). Por ejemplo, en $\mathbb{R}^3$, cualquier recta que pase por el origen es un subespacio de dimensión 1, y cualquier plano que pase por el origen es un subespacio de dimensión 2.

\section{Combinación lineal, independencia lineal, base y dimensión}

\subsection{Combinación lineal}
Dados vectores $\mathbf{v}_1, \dots, \mathbf{v}_k \in V$ y escalares $c_1, \dots, c_k \in \mathbb{R}$, la expresión
\[
c_1 \mathbf{v}_1 + c_2 \mathbf{v}_2 + \dots + c_k \mathbf{v}_k
\]
se llama \textbf{combinación lineal} de esos vectores.

\subsection{Independencia lineal}
Los vectores $\mathbf{v}_1, \dots, \mathbf{v}_k$ son \textbf{linealmente independientes} si la única combinación lineal que da el vector cero es aquella con todos los escalares cero:
\[
c_1 \mathbf{v}_1 + \dots + c_k \mathbf{v}_k = \mathbf{0} \quad \Rightarrow \quad c_1 = c_2 = \dots = c_k = 0
\]
Si existen escalares no todos nulos que cumplan la igualdad, son linealmente dependientes.

\subsection{Base}
Un conjunto de vectores $\{\mathbf{v}_1, \dots, \mathbf{v}_n\}$ es una \textbf{base} de $V$ si:
\begin{enumerate}
    \item Son linealmente independientes.
    \item Todo vector de $V$ puede escribirse como combinación lineal de ellos (es decir, generan $V$).
\end{enumerate}
La \textbf{dimensión} de $V$, denotada $\dim(V)$, es el número de elementos de una base (es único para un espacio dado).

Por ejemplo, la base canónica de $\mathbb{R}^3$ es $\{(1,0,0), (0,1,0), (0,0,1)\}$, y $\dim(\mathbb{R}^3)=3$.

\subsection{Coordenadas de un vector en una base}
Si $\mathcal{B} = \{\mathbf{v}_1, \dots, \mathbf{v}_n\}$ es una base, cualquier vector $\mathbf{v}$ se expresa de forma única como
\[
\mathbf{v} = \alpha_1 \mathbf{v}_1 + \dots + \alpha_n \mathbf{v}_n
\]
Los escalares $\alpha_i$ son las \textbf{coordenadas} de $\mathbf{v}$ en la base $\mathcal{B}$, y se representan como un vector columna $[\mathbf{v}]_{\mathcal{B}} = (\alpha_1, \dots, \alpha_n)^\top$.

\section{Cambio de base}

\subsection{Matriz de cambio de base}
Supongamos que tenemos dos bases de $V$: $\mathcal{B} = \{\mathbf{u}_1, \dots, \mathbf{u}_n\}$ y $\mathcal{C} = \{\mathbf{w}_1, \dots, \mathbf{w}_n\}. Queremos pasar las coordenadas de un vector en base $\mathcal{B}$ a coordenadas en base $\mathcal{C}$.

Sea $\mathbf{P}_{\mathcal{C} \leftarrow \mathcal{B}}$ la matriz cuyas columnas son las coordenadas de los vectores de $\mathcal{B}$ expresados en base $\mathcal{C}$. Entonces:
\[
[\mathbf{v}]_{\mathcal{C}} = \mathbf{P}_{\mathcal{C} \leftarrow \mathcal{B}} \; [\mathbf{v}]_{\mathcal{B}}
\]

Análogamente, la matriz inversa $\mathbf{P}_{\mathcal{B} \leftarrow \mathcal{C}} = \mathbf{P}_{\mathcal{C} \leftarrow \mathcal{B}}^{-1}$ transforma en la dirección contraria.

\subsection{Ejemplo en $\mathbb{R}^2$}
Considere la base canónica $\mathcal{C} = \{(1,0), (0,1)\}$ y otra base $\mathcal{B} = \{(1,1), (1,-1)\}$. Expresamos los vectores de $\mathcal{B}$ en base canónica:
\[
(1,1) = 1\cdot(1,0) + 1\cdot(0,1) \quad \Rightarrow \quad [ (1,1) ]_{\mathcal{C}} = \begin{pmatrix} 1 \\ 1 \end{pmatrix}
\]
\[
(1,-1) = 1\cdot(1,0) + (-1)\cdot(0,1) \quad \Rightarrow \quad [ (1,-1) ]_{\mathcal{C}} = \begin{pmatrix} 1 \\ -1 \end{pmatrix}
\]
Por tanto, la matriz de cambio de base de $\mathcal{B}$ a $\mathcal{C}$ es:
\[
\mathbf{P}_{\mathcal{C} \leftarrow \mathcal{B}} = \begin{pmatrix} 1 & 1 \\ 1 & -1 \end{pmatrix}
\]
Si un vector tiene coordenadas $[2,3]_{\mathcal{B}}$ (es decir, $2(1,1)+3(1,-1) = (5,-1)$), sus coordenadas en base canónica son:
\[
\begin{pmatrix} 1 & 1 \\ 1 & -1 \end{pmatrix} \begin{pmatrix} 2 \\ 3 \end{pmatrix} = \begin{pmatrix} 5 \\ -1 \end{pmatrix}
\]

\subsection{Aplicación en NLP: embeddings}
En procesamiento de lenguaje natural, representamos palabras mediante vectores densos (embeddings). Podemos pensar que cada palabra tiene coordenadas en una base de ``características semánticas''. Cambiar de base equivale a transformar la representación (por ejemplo, de one-hot a un espacio continuo). Las capas de embedding aprendidas son esencialmente matrices de cambio de base.

\section{Transformaciones lineales}

\subsection{Definición}
Una función $T: V \to W$ (donde $V$ y $W$ son espacios vectoriales) es una \textbf{transformación lineal} si cumple:
\begin{enumerate}
    \item $T(\mathbf{u} + \mathbf{v}) = T(\mathbf{u}) + T(\mathbf{v})$ para todos $\mathbf{u}, \mathbf{v} \in V$.
    \item $T(c\mathbf{v}) = c T(\mathbf{v})$ para todo escalar $c$ y $\mathbf{v} \in V$.
\end{enumerate}

\subsection{Representación matricial}
Si $\dim(V)=n$, $\dim(W)=m$, y elegimos bases $\mathcal{B}_V$ y $\mathcal{B}_W$, entonces $T$ puede representarse mediante una matriz $\mathbf{A} \in \mathbb{R}^{m \times n}$ tal que:
\[
[T(\mathbf{v})]_{\mathcal{B}_W} = \mathbf{A} \; [\mathbf{v}]_{\mathcal{B}_V}
\]
La columna $j$ de $\mathbf{A}$ contiene las coordenadas de $T(\mathbf{v}_j)$ en base $\mathcal{B}_W$, donde $\mathbf{v}_j$ es el $j$-ésimo vector de la base de $V$.

\subsection{Ejemplos geométricos}
\begin{itemize}
    \item \textbf{Rotación} en el plano (ángulo $\theta$):
    \[
    \begin{pmatrix} \cos\theta & -\sin\theta \\ \sin\theta & \cos\theta \end{pmatrix}
    \]
    \item \textbf{Escalado}:
    \[
    \begin{pmatrix} s_x & 0 \\ 0 & s_y \end{pmatrix}
    \]
    \item \textbf{Cizallamiento} (shear) horizontal:
    \[
    \begin{pmatrix} 1 & k \\ 0 & 1 \end{pmatrix}
    \]
    \item \textbf{Proyección} sobre el eje $x$:
    \[
    \begin{pmatrix} 1 & 0 \\ 0 & 0 \end{pmatrix}
    \]
\end{itemize}

\subsection{Composición de transformaciones}
La composición de transformaciones lineales corresponde al producto de sus matrices. Si $T: V \to W$ y $S: W \to U$, entonces $S \circ T$ tiene matriz $\mathbf{B}\mathbf{A}$ (siempre que las bases sean consistentes).

\section{Aplicaciones en Inteligencia Artificial}

\subsection{Deep Learning: capas fully connected}
Una capa totalmente conectada (linear layer) en una red neuronal realiza la transformación:
\[
\mathbf{y} = \mathbf{W}\mathbf{x} + \mathbf{b}
\]
La parte lineal $\mathbf{W}\mathbf{x}$ es una transformación lineal del espacio de entrada al espacio de salida. Cada columna de $\mathbf{W}$ representa cómo se transforma cada característica de entrada. Las redes profundas apilan varias de estas transformaciones intercaladas con no linealidades, deformando el espacio de entrada hasta que los datos se vuelven linealmente separables.

\begin{center}
\begin{tikzpicture}
\draw[->] (-1,0) -- (3,0) node[right] {Caract. 1};
\draw[->] (0,-1) -- (0,3) node[above] {Caract. 2};
\foreach \x/\y in {0.5/0.5, 1.5/2, 2/1, 0.8/1.8, 2.2/0.8} {
    \filldraw[blue] (\x,\y) circle (2pt);
}
\draw[->, red, thick] (0,0) -- (1,0.5);
\draw[->, red, thick] (0,0) -- (0.5,1);
\end{tikzpicture}
\captionof{figure}{Transformación lineal: los vectores base cambian.}
\end{center}

\subsection{NLP: embeddings y cambio de base}
En modelos como Word2Vec o GloVe, las palabras se representan mediante vectores en un espacio semántico. Las operaciones algebraicas (por ejemplo, vector('rey') - vector('hombre') + vector('mujer') ≈ vector('reina')) tienen sentido porque las relaciones semánticas corresponden a transformaciones lineales en ese espacio. Un cambio de base puede interpretarse como una reparametrización del espacio semántico.

\subsection{Visión por computador: data augmentation}
Para aumentar la variabilidad de los datos de entrenamiento, aplicamos transformaciones afines a las imágenes: rotaciones, escalados, traslaciones, cizallamientos. Estas son transformaciones lineales (o afines) que modifican las coordenadas de los píxeles. Por ejemplo, rotar una imagen equivale a aplicar la matriz de rotación a cada coordenada.

\section{Visualización con Python}

\subsection{Transformaciones en 2D con Matplotlib}
Ilustraremos cómo una matriz transforma un conjunto de puntos (un cuadrado) y cómo se ven los vectores base.

\begin{lstlisting}[language=Python]
import numpy as np
import matplotlib.pyplot as plt

def plot_transform(A, puntos, titulo):
    puntos_trans = A @ puntos
    plt.figure(figsize=(8,4))
    plt.subplot(1,2,1)
    plt.scatter(puntos[0,:], puntos[1,:], c='blue', alpha=0.5)
    plt.quiver(0,0, 1,0, angles='xy', scale_units='xy', scale=1, color='r', label='e1')
    plt.quiver(0,0, 0,1, angles='xy', scale_units='xy', scale=1, color='g', label='e2')
    plt.xlim(-3,3); plt.ylim(-3,3); plt.grid(); plt.axis('equal'); plt.title('Original')
    plt.legend()
    
    plt.subplot(1,2,2)
    plt.scatter(puntos_trans[0,:], puntos_trans[1,:], c='blue', alpha=0.5)
    e1_trans = A @ [1,0]
    e2_trans = A @ [0,1]
    plt.quiver(0,0, e1_trans[0], e1_trans[1], angles='xy', scale_units='xy', scale=1, color='r', label='e1\'')
    plt.quiver(0,0, e2_trans[0], e2_trans[1], angles='xy', scale_units='xy', scale=1, color='g', label='e2\'')
    plt.xlim(-3,3); plt.ylim(-3,3); plt.grid(); plt.axis('equal'); plt.title('Transformado')
    plt.legend()
    plt.suptitle(titulo)
    plt.tight_layout()
    plt.show()

# Crear un conjunto de puntos (un cuadrado)
x = np.linspace(-1,1,10)
y = np.linspace(-1,1,10)
X, Y = np.meshgrid(x,y)
puntos = np.vstack([X.ravel(), Y.ravel()])

# Rotación 30 grados
theta = np.radians(30)
A_rot = np.array([[np.cos(theta), -np.sin(theta)],
                  [np.sin(theta),  np.cos(theta)]])
plot_transform(A_rot, puntos, 'Rotación 30°')

# Escalado
A_scale = np.array([[2, 0], [0, 0.5]])
plot_transform(A_scale, puntos, 'Escalado (2 en x, 0.5 en y)')

# Cizallamiento horizontal
A_shear = np.array([[1, 1.5], [0, 1]])
plot_transform(A_shear, puntos, 'Cizallamiento horizontal')
\end{lstlisting}

\subsection{Implementar una capa lineal con NumPy}
Simularemos una capa fully connected con NumPy, que aplica una transformación lineal a un batch de datos.

\begin{lstlisting}[language=Python]
class LinearLayer:
    def __init__(self, in_features, out_features):
        # Inicialización aleatoria de pesos
        self.W = np.random.randn(out_features, in_features) * 0.01
        self.b = np.zeros((out_features, 1))
    
    def forward(self, X):
        # X: (in_features, batch_size)
        return self.W @ X + self.b

# Ejemplo: batch de 5 vectores de 3 dimensiones, capa de salida 2 dimensiones
batch = np.random.randn(3, 5)
layer = LinearLayer(3, 2)
salida = layer.forward(batch)
print("Entrada shape:", batch.shape)
print("Salida shape:", salida.shape)
print("Matriz W shape:", layer.W.shape)
\end{lstlisting}

\subsection{Cambio de base con NumPy}
Dada una base $\mathcal{B}$ representada por una matriz $P$ (cuyas columnas son los vectores base), pasar de coordenadas en base canónica a coordenadas en base $\mathcal{B}$ se hace resolviendo $P x = v$, es decir $x = P^{-1} v$.

\begin{lstlisting}[language=Python]
# Base B: vectores columna (1,1) y (1,-1)
P = np.array([[1, 1],
              [1, -1]])

# Vector en base canónica
v_canon = np.array([5, -1])

# Coordenadas en base B
v_B = np.linalg.solve(P, v_canon)
print("Coordenadas en B:", v_B)  # debe ser [2, 3]

# Verificación
print("Reconstrucción:", P @ v_B)
\end{lstlisting}

\section{Notación en papers científicos}

En artículos de machine learning, es común encontrar expresiones como:

\begin{itemize}
    \item Capa lineal: $\mathbf{h} = \mathbf{W}\mathbf{x} + \mathbf{b}$, con $\mathbf{W} \in \mathbb{R}^{d_{\text{out}} \times d_{\text{in}}}$.
    \item Cambio de base en embeddings: $\mathbf{e}_{\text{word}} = \mathbf{E} \mathbf{1}_{\text{word}}$, donde $\mathbf{E}$ es la matriz de embedding (cada columna es el embedding de una palabra en la base canónica).
    \item Transformaciones afines en data augmentation: $T_{\theta}(\mathbf{p}) = \mathbf{R}_{\theta} \mathbf{p} + \mathbf{t}$, con $\mathbf{R}_{\theta}$ matriz de rotación.
    \item Descomposición de una matriz de pesos como cambio de base: $\mathbf{W} = \mathbf{U} \mathbf{D} \mathbf{V}^\top$ (SVD), donde $\mathbf{U}$ y $\mathbf{V}$ son bases ortonormales.
\end{itemize}

También se habla de ``espacio latente'' (latent space) como un espacio vectorial donde las representaciones internas del modelo son linealmente transformables.

\section{Ejercicios propuestos}

\begin{enumerate}
    \item Demuestra que el conjunto de matrices $2\times2$ con la suma y producto por escalar habituales es un espacio vectorial. ¿Cuál es su dimensión?
    \item En $\mathbb{R}^3$, determina si los vectores $(1,2,3)$, $(4,5,6)$ y $(7,8,9)$ son linealmente independientes.
    \item Encuentra la matriz de la transformación lineal que refleja un vector respecto a la línea $y=x$ en el plano.
    \item Usando Python, genera 100 puntos aleatorios en $[-1,1]\times[-1,1]$ y aplícales una rotación de 60 grados, un escalado de $(2,0.5)$ y un cizallamiento de factor 0.8. Visualiza los puntos originales y transformados.
    \item Investiga cómo se representa una convolución en una red neuronal como una transformación lineal (pista: matriz de Toeplitz).
\end{enumerate}

\section{Resumen}

En esta sesión hemos explorado:
\begin{itemize}
    \item La estructura de espacio vectorial que subyace a los datos en IA.
    \item Los conceptos de base, dimensión y cambio de base, fundamentales para entender representaciones.
    \item Las transformaciones lineales y su representación matricial.
    \item Aplicaciones concretas: capas de redes neuronales, embeddings, data augmentation.
    \item Visualización e implementación en Python.
    \item La notación típica en papers.
\end{itemize}
Estos conceptos son el andamiaje sobre el que se construyen modelos más complejos. En las próximas sesiones profundizaremos en autovalores, descomposiciones y otras herramientas avanzadas.

\end{document}